\section{Примеры и количественное сравнение}\label{sec:examples}
\subsection{Параметры через зазор и большая общая константа}
Рассмотрим
\begin{equation}\label{eq:exp_family}
 f_A(x,y)=e^x+Ay,\qquad
 Q=[0,R]\times[0,1],\qquad R,A>0,\qquad D=\sqrt{R^2+1}.
\end{equation}
Функция выпукла, минимум достигается в $(0,0)$, причём
$\nabla f_A(0,0)=(1,A)$.
Линейный минимум всегда равен $(0,0)$, поэтому
\[
 G(x,y)=xe^x+Ay,\qquad
 \|\nabla^2 f_A(x,y)\|=e^x\le1+xe^x\le1+G(x,y).
\]
Действительно, $(1-x)e^x$ не возрастает при $x\ge0$ и равна единице
в нуле. Значит, для всех $R,A$ подходят $L_0=L_1=1$,
а наименьшая общая граница гессиана на $Q$ равна $L=e^R$.

\begin{theorem}[Два правила шага на одном семействе]\label{thm:separation}
Начнём оба метода из $(R,0)$ в задаче \eqref{eq:exp_family}.
Если продолжать итерации без более ранней остановки с положительным
допуском, логарифмическое правило \eqref{eq:log_step} достигает точного
минимума не более чем за
\begin{equation}\label{eq:exp_upper}
 N_{\log}\le\left\lceil\frac{R\sqrt{R^2+1}}{\log2}\right\rceil
\end{equation}
шагов. Для обычного короткого шага
$\alpha=\min\{1,G/(L\|d\|^2)\}$ с $L=e^R$
достижение ошибки по функции не выше $\varepsilon>0$ требует
\begin{equation}\label{eq:exp_lower}
 N_L\ge\frac{e^R/(1+\varepsilon)-1}{e-1}.
\end{equation}
Точки, шаги, значения функции, зазоры, оценки ошибки и число
вычислений не зависят от $A$.
\end{theorem}
\begin{proof}
Все точки имеют $y=0$. Для логарифмического правила первая
координата обновляется по формуле
\[
 x_{\mathrm{next}}=\max\{0,x-h(x)\},\qquad
 h(x)=\frac1D\log\left(1+\frac{D}{x+e^{-x}}\right).
\]
На $[0,R]$ верно $x+e^{-x}\le R+e^{-R}\le D$.
Последнее следует из $e^R-e^{-R}\ge2R$: производная левой части
не меньше двух.
Поэтому $h(x)\ge(\log2)/D$, что доказывает \eqref{eq:exp_upper}.

Для обычного короткого шага
$x_{\mathrm{next}}=x-h$, где $h=\min\{x,e^{x-R}\}\le1$.
Из выпуклости экспоненты на $[0,1]$ следует
\[
 e^{-(x-h)}-e^{-x}
 =e^{-x}(e^h-1)
 \le(e-1)e^{-x}h\le(e-1)e^{-R}.
\]
В начале $e^{-x}=e^{-R}$. Ошибка не выше $\varepsilon$ требует
$e^{-x}\ge1/(1+\varepsilon)$, откуда получаем \eqref{eq:exp_lower}.
На грани $y=0$ значения функции, зазоры, направления и формулы
шага не зависят от $A$. То же верно для проверки \eqref{eq:stop}.
\end{proof}

При фиксированной точности верхняя граница для нового шага
имеет порядок $R^2$. Нижняя граница для обычного шага с $L=e^R$
имеет порядок $e^R$. Значения $L_0,L_1$ влияют на доказанную
оценку работы метода. Сравнивать лишь числа с разными
размерностями недостаточно.

\paragraph{Границы сравнения.}
В задаче \eqref{eq:exp_family} линейный оракул в любой точке возвращает
глобальный минимизатор. Поэтому обычный шаг $\alpha_k=2/(k+2)$,
считая $k$ от нуля, достигает минимума за одно обновление; то же
верно для точного поиска по отрезку. Теорема~\ref{thm:separation}
разделяет два конкретных правила шага и не является нижней границей
сложности всех вариантов FW или доказательством превосходства над
адаптивными конкурентами. Изменение $A$ при старте $(R,0)$ не создаёт
новых траекторий: это одномерная проверка чувствительности к нормальной
компоненте градиента.

\subsection{Невыпуклая функция с ненулевым градиентом в минимуме}
На $Q=[0,1]$ рассмотрим
\begin{equation}\label{eq:nonconvex}
 f(x)=\int_0^x e^{-2t^2}\,dt.
\end{equation}
Минимизатор $x^*=0$, производная $f'(0)=1$, а $f''(x)<0$ при $x>0$.
Кроме того,
\[
 G(x)=xe^{-2x^2},\qquad
 |f''(x)|=4G(x),\qquad
 xf'(x)\ge e^{-2}x\ge e^{-2}f(x).
\]
Условие \eqref{eq:gap_smoothness} выполнено при $L_0=0,L_1=4$,
а \eqref{eq:quasar} выполнено при $\gamma=e^{-2}$.
Это невыпуклый пример заявленного класса.

Из $x_0=1$ первый логарифмический шаг даёт
$x_1=1-(\log2)/4\in(1/2,1)$.
Функция $xe^{-2x^2}$ убывает по $x$ на $(1/2,1]$,
поэтому $G(x_1)>G(x_0)$, хотя $f(x_1)<f(x_0)$.
Именно поэтому доказательство остановки использует монотонную
величину $f(x_k)-\ell_k$ и не предполагает монотонности $G(x_k)$.
Норма градиента при приближении к минимуму на множестве растёт
к единице.

Здесь линейный оракул тоже возвращает минимизатор $0$, а точный
поиск по отрезку выбирает полный шаг. Пример проверяет принадлежность
невыпуклой функции классу при $L_0=0$ и немонотонность $G$,
но не демонстрирует сложную численную задачу.
